\documentclass[aspectratio=169,11pt]{beamer}

% =========================================================
% Apresentacao/apostila:
% Subprogramas em VHDL: funcoes e procedimentos
% Baseada nos arquivos da aula:
% a009a_timer.vhd, a009b_timer_tb.vhd, a009c_timer_proc.vhd,
% a009d_decod.vhd, a009e_timer_proc_top.vhd,
% a009f_timer_proc_tb.vhd, a009g_lab_1.md, a009h_lab_2.md, a009j_vga.vhd
%
% O PDF subprogramas.pdf foi considerado apenas como apoio conceitual:
% subprogramas, funcoes, procedimentos, modos de parametros e classes
% de objetos.
% =========================================================

\usepackage[utf8]{inputenc}
\usepackage[T1]{fontenc}
\usepackage[brazil]{babel}
\usepackage{lmodern}
\usepackage{microtype}
\usepackage{listings}
\usepackage{xcolor}
\usepackage{booktabs}
\usepackage{array}
\usepackage{upquote}
\usepackage{tikz}
\usetikzlibrary{arrows.meta, positioning, fit}

\usetheme{Boadilla}
\usecolortheme{dolphin}
\setbeamertemplate{navigation symbols}{}
\setbeamertemplate{itemize items}[circle]
\setbeamertemplate{enumerate items}[default]

\definecolor{vhdlblue}{RGB}{26,74,138}
\definecolor{vhdlgreen}{RGB}{30,120,70}
\definecolor{vhdlgray}{RGB}{245,246,248}
\definecolor{vhdlred}{RGB}{150,45,45}

\lstdefinestyle{vhdlstyle}{
    language=VHDL,
    basicstyle=\ttfamily\scriptsize,
    keywordstyle=\color{vhdlblue}\bfseries,
    commentstyle=\color{vhdlgreen},
    stringstyle=\color{vhdlred},
    numbers=left,
    numberstyle=\tiny\color{gray},
    stepnumber=1,
    numbersep=6pt,
    backgroundcolor=\color{vhdlgray},
    frame=single,
    rulecolor=\color{gray!35},
    tabsize=4,
    showstringspaces=false,
    columns=fullflexible,
    keepspaces=true,
    breaklines=true,
    captionpos=b,
    xleftmargin=1.2em,
    framexleftmargin=1.0em,
    literate=
        {á}{{\'a}}1 {à}{{\`a}}1 {ã}{{\~a}}1 {â}{{\^a}}1
        {é}{{\'e}}1 {ê}{{\^e}}1
        {í}{{\'i}}1
        {ó}{{\'o}}1 {õ}{{\~o}}1 {ô}{{\^o}}1
        {ú}{{\'u}}1
        {ç}{{\c{c}}}1
        {Á}{{\'A}}1 {Ã}{{\~A}}1
        {É}{{\'E}}1
        {Í}{{\'I}}1
        {Ó}{{\'O}}1 {Õ}{{\~O}}1
        {Ú}{{\'U}}1
        {Ç}{{\c{C}}}1
}
\lstset{style=vhdlstyle}

\title[Subprogramas em VHDL]{Subprogramas em VHDL}
\subtitle{Funções e procedimentos aplicados a timers, displays e VGA}
\author{Aula prática de VHDL}
\institute{Base: arquivos \texttt{Aula009}}
\date{}

\begin{document}

\begin{frame}
    \titlepage
\end{frame}

\begin{frame}{Roteiro da aula}
    \begin{enumerate}
        \item Timer sem procedimento: onde aparece a repetição.
        \item Testbench do timer: clock, reset e simulação rápida.
        \item Timer com procedimento: reutilização da lógica de incremento.
        \item Função para display de sete segmentos.
        \item Integração no top-level para a DE10-Lite.
        \item Exemplo VGA: procedimentos para varredura e desenho.
        \item Labs e boas práticas.
    \end{enumerate}
\end{frame}

\begin{frame}{O que são subprogramas em VHDL?}
    Subprogramas são blocos de código nomeados que podem ser reutilizados dentro de uma arquitetura, pacote, processo ou outro escopo declarativo.

    \vspace{0.5em}

    Em VHDL, os dois tipos principais são:
    \begin{itemize}
        \item \textbf{Função}: calcula e retorna um único valor.
        \item \textbf{Procedimento}: executa uma sequência de comandos e pode alterar sinais, variáveis ou produzir múltiplos resultados.
    \end{itemize}

    \vspace{0.5em}
    Nesta aula, a ideia não é usar subprogramas por estilo, mas porque os exemplos mostram uma necessidade real: evitar repetição e deixar explícita a intenção do hardware descrito.
\end{frame}

\begin{frame}{Função ou procedimento?}
    \begin{table}
        \centering
        \small
        \renewcommand{\arraystretch}{1.35}
        \setlength{\tabcolsep}{8pt}
        \begin{tabular}{>{\bfseries}p{0.22\textwidth} p{0.29\textwidth} p{0.38\textwidth}}
            \toprule
            Critério & Função & Procedimento \\
            \midrule
            Resultado & Retorna um valor & Pode ter várias saídas \\
            Uso típico & Expressões e conversões & Sequências reutilizáveis de comandos \\
            Efeito esperado & Sem efeito colateral & Pode alterar sinais ou variáveis \\
            Exemplo da aula & Decodificar um dígito & Incrementar contador com estouro \\
            \bottomrule
        \end{tabular}
    \end{table}

    \vspace{0.5em}
    Regra prática: se a operação parece uma expressão, use função. Se parece uma pequena rotina de controle, use procedimento.

    % Nota de fala: enfatizar que essa regra é didática. Em projetos maiores, a escolha também depende de estilo, pacotes e política de síntese.
\end{frame}

\begin{frame}[fragile]{Classes de objetos em parâmetros}
    Em subprogramas, o parâmetro formal pode indicar a classe do objeto recebido. Funções normalmente usam entradas; procedimentos podem também alterar sinais ou variáveis.

\begin{columns}[T,onlytextwidth]
\begin{column}{0.47\textwidth}
\textbf{Função}
\begin{lstlisting}
function decodificando(
    valor : integer
) return std_logic_vector is
begin
    -- calcula e retorna um valor
end function;
\end{lstlisting}
\end{column}

\begin{column}{0.50\textwidth}
\textbf{Procedimento}
\begin{lstlisting}
procedure incrementa(
    signal   contador : inout unsigned;
    constant max      : in    positive;
    constant habilita : in    boolean;
    variable estouro  : out   boolean
) is
\end{lstlisting}
\end{column}
\end{columns}

    \begin{itemize}
        \item Em função, \texttt{valor : integer} significa \texttt{constant valor : in integer}; normalmente omitimos o padrão.
        \item \texttt{variable} não pode ser usada como classe de parâmetro formal em funções.
        \item \texttt{signal}: conectado a um sinal do circuito; atribuições usam \texttt{<=}.
        \item Em procedimentos, \texttt{variable} é útil para resultado imediato dentro de um processo; atribuições usam \texttt{:=}.
        \item \texttt{file}: existe em VHDL, mas é voltado a simulação e não é usado nestes exemplos.
    \end{itemize}
\end{frame}

\begin{frame}{Funções e procedimentos: modos}
    Além da classe do objeto, o parâmetro tem um modo:

    \begin{itemize}
        \item \texttt{in}: o subprograma lê o valor.
        \item \texttt{out}: o subprograma produz um valor.
        \item \texttt{inout}: o subprograma lê o valor atual e também escreve um novo valor.
    \end{itemize}

    \vspace{0.5em}
    Em uma \textbf{função}, os parâmetros são entradas e a saída principal é o \texttt{return}. Por isso ela combina bem com expressões, como o decodificador de sete segmentos.

    \vspace{0.5em}
    Em um \textbf{procedimento}, \texttt{out} e \texttt{inout} são naturais: é o caso de \texttt{incrementa}, que altera o contador e devolve a flag de estouro.
\end{frame}

\begin{frame}[fragile]{Chamar dentro ou fora de um process}
    Dentro de um \texttt{process}, a chamada participa da sequência daquele processo. No timer, isso é essencial porque as chamadas acontecem na borda de clock.

\begin{lstlisting}
elsif falling_edge(clk) then
    incrementa(ticks,  clockFreq, true, controle);
    incrementa(segReg, 60, controle, controle);
    incrementa(minReg, 60, controle, controle);
    incrementa(horReg, 24, controle, controle);
end if;
\end{lstlisting}

    Fora de um \texttt{process}, funções aparecem naturalmente em atribuições concorrentes:

\begin{lstlisting}
d0 <= decodificando(unidade);
\end{lstlisting}

    \vspace{0.2em}
    Procedimentos também podem ter chamada concorrente em VHDL, mas em projetos didáticos e sintetizáveis costuma ser mais claro chamá-los dentro de processos.
\end{frame}

\section{Timer sem procedimento}

\begin{frame}[fragile]{Exemplo inicial: timer sem procedimento}
    O arquivo \texttt{a009a\_timer.vhd} descreve um timer com um contador interno de ticks e três contadores visíveis na interface.

\begin{lstlisting}
port(
    clk  : in  std_logic;
    nRst : in  std_logic;
    s    : out unsigned(5 downto 0);
    m    : out unsigned(5 downto 0);
    h    : out unsigned(4 downto 0)
);
\end{lstlisting}

\begin{lstlisting}
architecture rtl of a009a_timer is
    signal ticks  : unsigned(bits_necessarios(clockFreq) - 1 downto 0) := (others => '0');
    signal segReg : unsigned(5 downto 0) := (others => '0');
    -- minReg e horReg seguem a mesma ideia
begin
    -- logica sequencial do timer
\end{lstlisting}

    \vspace{0.2em}
    \texttt{ticks} pertence à arquitetura: divide o clock até formar 1 segundo. Já \texttt{s}, \texttt{m} e \texttt{h} são portas \texttt{out}, conectáveis ao testbench, ao top-level ou aos displays.
\end{frame}

\begin{frame}[fragile]{Reset e contador de ticks}
    O processo primeiro trata o reset assíncrono ativo em nível baixo. Depois, a cada borda de descida, \texttt{ticks} transforma ciclos de clock em segundos.

\begin{lstlisting}
if nRst = '0' then
    ticks  <= (others => '0');
    segReg <= (others => '0');
    minReg <= (others => '0');
    horReg <= (others => '0');

elsif falling_edge(clk) then
    if ticks = to_unsigned(clockFreq - 1, ticks'length) then
        ticks <= (others => '0');
        -- aqui entram segundos, minutos e horas
    else
        ticks <= ticks + 1;
    end if;
end if;
\end{lstlisting}

    \vspace{0.5em}
    Em simulação, \texttt{clockFreq} pode ser pequeno, como 10 Hz. No hardware real da DE10-Lite, o clock principal é de 50 MHz.
\end{frame}

\begin{frame}[fragile]{Incremento em cascata}
    Quando \texttt{ticks} estoura, o timer incrementa segundos. Se segundos estouram, incrementa minutos. Se minutos estouram, incrementa horas.

\begin{lstlisting}
if segReg = to_unsigned(59, segReg'length) then
    segReg <= (others => '0');

    if minReg = to_unsigned(59, minReg'length) then
        minReg <= (others => '0');

        if horReg = to_unsigned(23, horReg'length) then
            horReg <= (others => '0');
        else
            horReg <= horReg + 1;
        end if;
    else
        minReg <= minReg + 1;
    end if;
else
    segReg <= segReg + 1;
end if;
\end{lstlisting}

    \vspace{0.4em}
    A estrutura funciona, mas o padrão ``incrementa ou volta para zero'' aparece várias vezes.
\end{frame}

\begin{frame}{O que este código nos ensina?}
    O timer sem procedimento é um bom primeiro passo porque deixa a lógica visível:
    \begin{itemize}
        \item existe um contador base;
        \item cada contador possui um limite;
        \item ao chegar ao limite, ele volta para zero;
        \item o estouro de um contador habilita o próximo.
    \end{itemize}

    \vspace{0.5em}
    A motivação para o procedimento nasce daqui: a lógica de incremento com estouro é a mesma para \texttt{ticks}, \texttt{s}, \texttt{m} e \texttt{h}; mudam apenas o limite e a condição de habilitação.

    % Nota de fala: pedir aos alunos para identificarem "o algoritmo repetido" antes de mostrar o procedimento.
\end{frame}


\section{Testbench do timer}

\begin{frame}[fragile]{Testbench do timer}
    O arquivo \texttt{a009b\_timer\_tb.vhd} instancia o timer original e reduz a frequência para tornar a simulação curta.

\begin{lstlisting}
constant clockFreq : positive := 10;
constant clkPer    : time := 1000 ms / clockFreq;

signal clk, nRst : std_logic := '0';
signal segundo   : unsigned(5 downto 0);
signal minuto    : unsigned(5 downto 0);
signal hora      : unsigned(4 downto 0);
\end{lstlisting}

    \vspace{0.4em}
    Em vez de esperar 50 milhões de ciclos para ver um segundo passar, o testbench usa 10 ciclos por segundo simulado e verifica automaticamente o primeiro transbordo de segundos para minutos.
\end{frame}

\begin{frame}[fragile]{Instanciação do DUT}
    O DUT é o \textit{device under test}: o circuito que queremos exercitar na simulação.

\begin{lstlisting}
i_timer: entity work.a009a_timer(rtl)
    generic map (
        clockFreq => clockFreq
    )
    port map (
        clk  => clk,
        nRst => nRst,
        s    => segundo,
        m    => minuto,
        h    => hora
    );
\end{lstlisting}

    \vspace{0.4em}
    O \texttt{generic map} troca o valor padrão do projeto por um valor adequado à simulação.
\end{frame}

\begin{frame}[fragile]{Clock e reset no testbench}
    O clock é gerado por uma atribuição concorrente. O reset começa ativo e depois é liberado.

\begin{lstlisting}
clk <= not clk after clkPer / 2;

process is
begin
    wait until rising_edge(clk);
    wait until rising_edge(clk);

    nRst <= '1';

    wait;
end process;
\end{lstlisting}

    \vspace{0.4em}
    Mesmo que o DUT use \texttt{falling\_edge(clk)}, o testbench aguarda duas bordas de subida apenas para dar um intervalo inicial antes de liberar o reset.
\end{frame}

\section{Timer com procedimento}

\begin{frame}[fragile]{Timer com procedimento}
    O arquivo \texttt{a009c\_timer\_proc.vhd} preserva a ideia do timer, mas encapsula o padrão de incremento.

\begin{lstlisting}
procedure incrementa(
    signal   contador : inout unsigned;
    constant max      : in    positive;
    constant habilita : in    boolean;
    variable estouro  : out   boolean
) is
begin
    estouro := false;
    if habilita then
        if contador = to_unsigned(max - 1, contador'length) then
            contador <= to_unsigned(0, contador'length);
            estouro := true;
        else
            contador <= contador + 1;
        end if;
    end if;
end procedure;
\end{lstlisting}

    \vspace{0.3em}
    Este procedimento descreve o comportamento: incrementar se estiver habilitado e avisar quando houver estouro.
\end{frame}

\begin{frame}{Por que estes parâmetros?}
    \begin{itemize}
        \item \texttt{signal contador : inout unsigned}: o procedimento altera um registrador do circuito, como \texttt{ticks}, \texttt{segReg}, \texttt{minReg} ou \texttt{horReg}.
        \item \texttt{constant max : in positive}: o limite é apenas lido; não deve ser modificado.
        \item \texttt{constant habilita : in boolean}: a condição de avanço também é apenas lida.
        \item \texttt{variable estouro : out boolean}: a resposta é usada imediatamente dentro do processo para habilitar o próximo contador.
    \end{itemize}

    \vspace{0.5em}
    A classe do objeto importa: sinais modelam conexões registradas ou combinacionais; variáveis dentro do processo atualizam imediatamente no fluxo sequencial.
\end{frame}

\begin{frame}[fragile]{Chamadas em cascata}
    No processo principal, a mesma variável \texttt{controle} carrega o estouro de um contador para o próximo.

\begin{lstlisting}
process(clk, nRst) is
    variable controle : boolean;
begin
    if nRst = '0' then
        ticks  <= (others => '0');
        segReg <= (others => '0');
        minReg <= (others => '0');
        horReg <= (others => '0');

    elsif falling_edge(clk) then
        incrementa(ticks,  clockFreq, true, controle);
        incrementa(segReg, 60, controle, controle);
        incrementa(minReg, 60, controle, controle);
        incrementa(horReg, 24, controle, controle);
    end if;
end process;
\end{lstlisting}

    \vspace{0.3em}
    A leitura é quase uma frase: incrementa ticks sempre; se estourar, incrementa segundos; se estourar, incrementa minutos; se estourar, incrementa horas.
\end{frame}

\begin{frame}{Comparação: antes e depois}
    \begin{table}
        \centering
        \small
        \begin{tabular}{p{0.42\textwidth} p{0.42\textwidth}}
            \toprule
            \textbf{Timer sem procedimento} & \textbf{Timer com procedimento} \\
            \midrule
            Mostra toda a lógica explicitamente. & Dá nome à lógica repetida. \\
            Usa vários \texttt{if} aninhados. & Usa chamadas sequenciais em cascata. \\
            É fácil de acompanhar no primeiro contato. & É mais fácil de manter e ampliar. \\
            A repetição aparece no corpo do processo. & A repetição fica concentrada em \texttt{incrementa}. \\
            \bottomrule
        \end{tabular}
    \end{table}

    \vspace{0.5em}
    O procedimento não cria uma ``função de software em execução''. Ele é uma forma de escrever a descrição do hardware de maneira organizada.
\end{frame}

\begin{frame}{Observações de síntese e estilo}
    \begin{itemize}
        \item O procedimento é sintetizável porque contém lógica compatível com hardware: comparações, atribuições e controle simples.
        \item Com \texttt{unsigned}, a largura do contador fica explícita e as comparações usam \texttt{to\_unsigned}.
        \item As portas são \texttt{out}; os valores são guardados em sinais internos como \texttt{segReg}, \texttt{minReg} e \texttt{horReg}.
        \item Usar subprogramas não remove a necessidade de pensar em registradores, bordas de clock e reset.
    \end{itemize}

    % Nota de fala: reforçar que subprograma em HDL organiza a descrição; não significa execução dinâmica como uma chamada de função em C.
\end{frame}

\section{Testbench do timer com procedimento}

\begin{frame}[fragile]{Testbench do timer com procedimento}
    O arquivo \texttt{a009f\_timer\_proc\_tb.vhd} repete a estratégia do testbench anterior.

\begin{lstlisting}
constant clockFreq : positive := 10;
constant clkPer    : time := 1000 ms / clockFreq;

signal clk, nRst : std_logic := '0';
signal segundo   : unsigned(5 downto 0);
signal minuto    : unsigned(5 downto 0);
signal hora      : unsigned(4 downto 0);
\end{lstlisting}

    \vspace{0.4em}
    A interface observada é a mesma em intenção: clock, reset, segundos, minutos e horas. Agora os contadores aparecem como vetores \texttt{unsigned}.
\end{frame}

\begin{frame}[fragile]{Mesmo teste, outra implementação}
\begin{lstlisting}
i_timer : entity work.a009c_timer_proc(rtl)
    generic map(
        clockFreq => clockFreq
    )
    port map(
        clk  => clk,
        nRst => nRst,
        s    => segundo,
        m    => minuto,
        h    => hora
    );
\end{lstlisting}

    \vspace{0.5em}
    Isso é uma ideia importante em projeto digital: podemos refatorar a implementação interna mantendo a interface. O testbench ajuda a verificar se o comportamento externo continua coerente.
\end{frame}

\section{Função para display}

\begin{frame}[fragile]{Função para display de sete segmentos}
    O arquivo \texttt{a009d\_decod.vhd} usa uma função para converter um valor decimal em padrão de segmentos.

\begin{lstlisting}
function decodificando(valor: integer)
    return std_logic_vector is

    variable saida : std_logic_vector(0 to 6);
begin
    case valor is
        when 0 => saida := not "1111110";
        when 1 => saida := not "0110000";
        when 2 => saida := not "1101101";
        -- ...
        when 9 => saida := not "1111011";
        when others => saida := "0000000";
    end case;

    return saida;
end decodificando;
\end{lstlisting}

    \vspace{0.3em}
    Aqui a função é natural: entra um inteiro, sai um vetor de sete bits.
\end{frame}

\begin{frame}{Por que função é adequada aqui?}
    A decodificação do display tem uma característica simples:
    \begin{itemize}
        \item não precisa alterar sinais por conta própria;
        \item não precisa devolver vários resultados;
        \item representa uma tabela de conversão;
        \item pode ser usada diretamente em atribuições.
    \end{itemize}

    \vspace{0.5em}
    A função \texttt{decodificando} transforma uma decisão combinacional em um bloco nomeado. Isso melhora a leitura do processo que separa dezena e unidade.
\end{frame}

\begin{frame}[fragile]{Decomposição em dezena e unidade}
    O processo do decodificador separa o número decimal em dois dígitos.

\begin{lstlisting}
process(contador) is
    variable valor   : integer range 0 to 63;
    variable dezena  : integer range 0 to 6;
    variable unidade : integer range 0 to 9;
begin
    valor   := to_integer(contador);
    dezena  := valor / 10;
    unidade := valor mod 10;

    d1 <= decodificando(dezena);
    d0 <= decodificando(unidade);
end process;
\end{lstlisting}

    \vspace{0.4em}
    \texttt{to\_integer} converte o \texttt{unsigned} para cálculo decimal. Depois, \texttt{/} faz divisão inteira e \texttt{mod} calcula o resto.
\end{frame}

\begin{frame}{Detalhe do hardware da placa}
    Os padrões são escritos com \texttt{not "1111110"}, \texttt{not "0110000"} e assim por diante.

    \vspace{0.5em}
    Isso indica uma característica comum dos displays de sete segmentos em kits FPGA: os segmentos podem ser ativos em nível baixo. Nesse caso:
    \begin{itemize}
        \item bit em \texttt{'0'} acende o segmento;
        \item bit em \texttt{'1'} apaga o segmento;
        \item o \texttt{not} adapta uma tabela pensada como ativo alto para a saída ativa baixa.
    \end{itemize}
\end{frame}

\section{Top-level}

\begin{frame}[fragile]{Integração no top-level}
    O arquivo \texttt{a009e\_timer\_proc\_top.vhd} combina o timer com três instâncias do decodificador.

\begin{lstlisting}
entity a009e_timer_proc_top is
    generic(
        clockFreq : positive := 50e6
    );
    port(
        MAX10_CLK1_50 : in  std_logic;
        KEY           : in  std_logic_vector(0 downto 0);
        HEX0, HEX1    : out std_logic_vector(0 to 6);
        HEX2, HEX3    : out std_logic_vector(0 to 6);
        HEX4, HEX5    : out std_logic_vector(0 to 6)
    );
end entity;
\end{lstlisting}

    \vspace{0.3em}
    Esta entidade é a ponte entre o projeto VHDL e os pinos do kit DE10-Lite.
\end{frame}

\begin{frame}{Integração do timer no top-level}
    \centering
    \begin{tikzpicture}[
        font=\scriptsize,
        block/.style={draw, rounded corners=2pt, thick, align=center, minimum height=0.82cm, fill=blue!6},
        io/.style={draw, rounded corners=2pt, thick, align=center, minimum height=0.62cm, fill=gray!10},
        wire/.style={-{Latex[length=2mm]}, thick},
        bus/.style={-{Latex[length=2mm]}, very thick},
        note/.style={align=center, font=\tiny}
    ]
        \node[io, minimum width=1.8cm] (clk) {MAX10\_CLK1\_50\\clock 50 MHz};
        \node[io, below=0.55cm of clk, minimum width=1.8cm] (key) {KEY(0)\\reset};

        \node[block, right=1.4cm of clk, yshift=-0.28cm, minimum width=2.2cm, minimum height=1.65cm] (timer) {\texttt{a009c}\\timer com\\procedimento};

        \node[block, right=1.75cm of timer, yshift=0.95cm, minimum width=2.15cm] (decs) {\texttt{a009d}\\decod seg.};
        \node[block, right=1.75cm of timer, minimum width=2.15cm] (decm) {\texttt{a009d}\\decod min.};
        \node[block, right=1.75cm of timer, yshift=-0.95cm, minimum width=2.15cm] (dech) {\texttt{a009d}\\decod hora};

        \node[io, right=1.4cm of decs, minimum width=1.55cm] (hexs) {HEX0\\HEX1};
        \node[io, right=1.4cm of decm, minimum width=1.55cm] (hexm) {HEX2\\HEX3};
        \node[io, right=1.4cm of dech, minimum width=1.55cm] (hexh) {HEX4\\HEX5};

        \draw[wire] (clk.east) -- node[above, note] {\texttt{clk}} (timer.west);
        \draw[wire] (key.east) -- node[below, note] {\texttt{nRst}} (timer.west);

        \draw[bus] (timer.east) -- ++(0.65,0) |- node[pos=0.25, above, note] {\texttt{segundos}} (decs.west);
        \draw[bus] (timer.east) -- node[above, note] {\texttt{minutos}} (decm.west);
        \draw[bus] (timer.east) -- ++(0.65,0) |- node[pos=0.25, below, note] {\texttt{horas}} (dech.west);

        \draw[bus] (decs.east) -- node[above, note] {7 segmentos} (hexs.west);
        \draw[bus] (decm.east) -- node[above, note] {7 segmentos} (hexm.west);
        \draw[bus] (dech.east) -- node[above, note] {7 segmentos} (hexh.west);

        \node[draw, dashed, rounded corners=3pt, fit=(timer)(decs)(decm)(dech), inner sep=0.25cm, label={[font=\scriptsize]above:\texttt{a009e\_timer\_proc\_top}}] {};
    \end{tikzpicture}

    \vspace{0.6em}
    O top-level não calcula o tempo nem decodifica segmentos: ele organiza as conexões entre os blocos.
\end{frame}

\begin{frame}[fragile]{Três decodificadores}
    Cada contador do timer alimenta um par de displays.

\begin{lstlisting}
decod_s: entity work.a009d_decod(rtl)
    port map(contador => segundos, d0 => HEX0, d1 => HEX1);

decod_m: entity work.a009d_decod(rtl)
    port map(contador => minutos, d0 => HEX2, d1 => HEX3);

decod_h: entity work.a009d_decod(rtl)
    port map(contador => resize(horas, 6), d0 => HEX4, d1 => HEX5);
\end{lstlisting}

    \vspace{0.4em}
    A reutilização aparece em dois níveis: a função é reutilizada dentro do decodificador, e o módulo decodificador é reutilizado três vezes no top-level.
\end{frame}

\begin{frame}[fragile]{Timer conectado à DE10-Lite}
\begin{lstlisting}
timer_0: entity work.a009c_timer_proc(rtl)
    generic map(
        clockFreq => clockFreq
    )
    port map(
        clk  => MAX10_CLK1_50,
        nRst => KEY(0),
        s    => segundos,
        m    => minutos,
        h    => horas
    );
\end{lstlisting}

    \vspace{0.4em}
    O projeto completo usa o clock de 50 MHz da placa, o botão \texttt{KEY(0)} como reset ativo baixo e os displays \texttt{HEX0} a \texttt{HEX5} para mostrar segundos, minutos e horas.
\end{frame}

\begin{frame}[fragile]{Testbench do top-level}
    O arquivo \texttt{a009e\_timer\_proc\_top\_tb.vhd} testa a integração completa com frequência reduzida.

\begin{lstlisting}
dut: entity work.a009e_timer_proc_top(main)
    generic map(clockFreq => 10)
    port map(
        MAX10_CLK1_50 => clk,
        KEY           => key,
        HEX0 => HEX0, HEX1 => HEX1,
        HEX2 => HEX2, HEX3 => HEX3,
        HEX4 => HEX4, HEX5 => HEX5
    );
\end{lstlisting}

    \vspace{0.3em}
    O teste verifica padrões nos displays: primeiro \texttt{00:00:01}, depois \texttt{00:01:00}. Assim, ele cobre timer, decodificadores e conexões do top-level.
\end{frame}

\begin{frame}{O que o top-level mostra sobre organização?}
    O top-level não deveria conter todos os detalhes internos do timer nem todos os detalhes da tabela de sete segmentos.

    \vspace{0.5em}
    Seu papel é conectar blocos:
    \begin{itemize}
        \item o timer produz valores numéricos;
        \item os decodificadores transformam valores em segmentos;
        \item os pinos da placa recebem clock, reset e sinais de display.
    \end{itemize}

    \vspace{0.5em}
    Subprogramas ajudam dentro dos módulos. Instanciação de entidades ajuda entre módulos. As duas ideias se complementam.
\end{frame}

\section{VGA}

\begin{frame}[fragile]{Exemplo VGA}
    O arquivo \texttt{a009j\_vga.vhd} gera três quadrados coloridos em um monitor VGA 800x600 a partir do clock de 50 MHz da DE10-Lite.

\begin{lstlisting}
entity a009j_vga is
    port (
        MAX10_CLK1_50  : in  std_logic;
        VGA_R          : out std_logic_vector(3 downto 0);
        VGA_G          : out std_logic_vector(3 downto 0);
        VGA_B          : out std_logic_vector(3 downto 0);
        VGA_VS, VGA_HS : out std_logic
    );
end entity;
\end{lstlisting}

    \vspace{0.3em}
    Este exemplo usa procedimentos para separar duas preocupações: varredura da tela e desenho.
\end{frame}

\begin{frame}[fragile]{Constantes de temporização}
    As constantes definem os intervalos horizontal e vertical: área visível, \textit{front porch}, pulso de sincronismo, \textit{back porch} e total.

\begin{lstlisting}
constant H_ACTIVE_VIDEO : integer := 800;
constant H_FRONT_PORCH  : integer := H_ACTIVE_VIDEO + 56;
constant H_SYNC         : integer := H_FRONT_PORCH + 120;
constant H_BACK_PORCH   : integer := H_SYNC + 64;
constant H_PIXELS       : integer := 1040;

constant V_ACTIVE_VIDEO : integer := 600;
constant V_FRONT_PORCH  : integer := V_ACTIVE_VIDEO + 37;
constant V_SYNC         : integer := V_FRONT_PORCH + 6;
constant V_BACK_PORCH   : integer := V_SYNC + 23;
constant V_PIXELS       : integer := 666;
\end{lstlisting}
\end{frame}

\begin{frame}[fragile]{Procedimento incrementa reutilizado}
    O mesmo padrão do timer reaparece no controlador VGA: incrementar um contador e informar estouro.

\begin{lstlisting}
incrementa(horizontal, H_PIXELS, true, controle);
incrementa(vertical, V_PIXELS, controle, controle);
\end{lstlisting}

    \vspace{0.5em}
    A diferença é o significado dos contadores:
    \begin{itemize}
        \item \texttt{horizontal}: posição do pixel dentro da linha.
        \item \texttt{vertical}: posição da linha dentro do quadro.
    \end{itemize}

    \vspace{0.4em}
    Quando o contador horizontal estoura, o vertical avança uma linha.
\end{frame}

\begin{frame}[fragile]{Sincronismos e área ativa}
    O processo de varredura calcula os pulsos de sincronismo e a região visível.

\begin{lstlisting}
if horizontal >= H_FRONT_PORCH and horizontal < H_SYNC then
    VGA_HS <= '0';
end if;

if vertical >= V_FRONT_PORCH and vertical < V_SYNC then
    VGA_VS <= '0';
end if;

if horizontal < H_ACTIVE_VIDEO and vertical < V_ACTIVE_VIDEO then
    active <= true;
end if;
\end{lstlisting}

    \vspace{0.4em}
    Esta parte é lógica de controle da tela. Ela não decide o que será desenhado; apenas informa onde o feixe lógico está.
\end{frame}

\begin{frame}[fragile]{Procedimento quadrado}
    O procedimento \texttt{quadrado} concentra a decisão geométrica de desenho.

\begin{lstlisting}
procedure quadrado(
    constant enable : in boolean;
    constant h, v   : in integer;
    constant x, y   : in integer;
    constant w      : in integer;
    constant color  : in std_logic_vector(11 downto 0);
    variable draw   : inout std_logic_vector(11 downto 0)
) is
begin
    if enable then
        if (v >= y - w/2 and v < y + w/2) and
           (h >= x - w/2 and h < x + w/2) then
            draw := color;
        end if;
    end if;
end procedure;
\end{lstlisting}

    \vspace{0.2em}
    O parâmetro \texttt{draw} é \texttt{variable inout} porque a cor pode ser modificada imediatamente dentro do processo de desenho.
\end{frame}

\begin{frame}[fragile]{Desenho com chamadas sucessivas}
    No processo de desenho, o fundo começa preto e cada chamada pode substituir a cor do pixel atual.

\begin{lstlisting}
draw := (others => '0');

quadrado(active, horizontal, vertical, 350, 275, 60, x"F00", draw);
quadrado(active, horizontal, vertical, 400, 300, 60, x"0F0", draw);
quadrado(active, horizontal, vertical, 450, 325, 60, x"00F", draw);

VGA_R <= draw(11 downto 8);
VGA_G <= draw(7 downto 4);
VGA_B <= draw(3 downto 0);
\end{lstlisting}

    \vspace{0.4em}
    A leitura fica direta: desenhe um quadrado vermelho, um verde e um azul, cada um com centro, tamanho e cor próprios.
\end{frame}

\begin{frame}{Controle de varredura versus lógica de desenho}
    O exemplo VGA separa bem duas tarefas:
    \begin{itemize}
        \item \textbf{Varredura}: contar pixels e linhas, gerar \texttt{VGA\_HS}, \texttt{VGA\_VS} e \texttt{active}.
        \item \textbf{Desenho}: decidir a cor RGB para a coordenada atual.
    \end{itemize}

    \vspace{0.5em}
    Os procedimentos ajudam porque dão nome às operações repetidas: \texttt{incrementa} para contadores e \texttt{quadrado} para teste geométrico de região.

    \vspace{0.5em}
    Observação didática: como \texttt{active} é um sinal atribuído no processo clockado, seu valor é atualizado após o ciclo de simulação atual. Em projetos de vídeo mais rigorosos, é comum alinhar sinais de controle e pixels com cuidado.
\end{frame}

\section{Labs}

\begin{frame}{Atividades de laboratório}
    As atividades de laboratório estão no repositório da disciplina, nos arquivos Markdown da aula:
    \begin{itemize}
        \item \texttt{a009g\_lab\_1.md}: conversor Celsius para Fahrenheit usando função.
        \item \texttt{a009h\_lab\_2.md}: gerador PWM usando procedimento.
    \end{itemize}

    \vspace{0.5em}
    A sequência dos labs acompanha a aula: primeiro uma função com retorno único, depois um procedimento para uma rotina de controle reutilizável.
\end{frame}

\begin{frame}{Como pensar os labs}
    \begin{itemize}
        \item No conversor de temperatura, a função é adequada porque recebe Celsius e retorna Fahrenheit.
        \item No PWM, o procedimento é adequado porque descreve uma sequência de controle baseada em período, contador e duty cycle.
        \item Os dois labs devem começar com constantes reduzidas em simulação antes de usar valores adequados ao clock da placa.
    \end{itemize}

    \vspace{0.5em}
    Antes de levar para a placa, simule uma versão pequena. Reduzir constantes em testbench é uma estratégia legítima para observar comportamento sem esperar milhões de ciclos.
\end{frame}

\section{Fechamento}

\begin{frame}{Resumo: função}
    Use função quando:
    \begin{itemize}
        \item a operação retorna um único valor;
        \item o código pode ser lido como uma expressão;
        \item não há necessidade de alterar sinais externos;
        \item a intenção é conversão, cálculo combinacional ou seleção de valor.
    \end{itemize}

    \vspace{0.5em}
    Exemplo da aula: \texttt{decodificando(valor)} retorna o padrão de sete segmentos para um dígito.
\end{frame}

\begin{frame}{Resumo: procedimento}
    Use procedimento quando:
    \begin{itemize}
        \item há uma sequência de comandos reutilizável;
        \item a operação altera um sinal ou variável passado como parâmetro;
        \item há mais de uma informação de saída;
        \item a lógica tem caráter de ação: incrementar, testar estouro, desenhar, atualizar.
    \end{itemize}

    \vspace{0.5em}
    Exemplos da aula: \texttt{incrementa} no timer e no VGA; \texttt{quadrado} no processo de desenho.
\end{frame}

\begin{frame}{Boas práticas}
    \begin{itemize}
        \item Comece com subprogramas simples e próximos do problema real.
        \item Evite subprogramas genéricos demais no início; eles podem esconder a lógica que o aluno ainda precisa enxergar.
        \item Dê nomes que expressem comportamento: \texttt{incrementa}, \texttt{decodificando}, \texttt{quadrado}.
        \item Teste com frequências e limites reduzidos antes de usar os valores reais do hardware.
        \item Para síntese, prefira faixas explícitas quando usar inteiros.
        \item Mantenha a interface do módulo estável quando refatorar a implementação interna.
    \end{itemize}
\end{frame}

\begin{frame}{Subprogramas e síntese}
    Em software, uma função costuma ser imaginada como uma chamada executada no tempo.

    \vspace{0.5em}
    Em VHDL sintetizável, subprogramas são principalmente uma forma de escrever a descrição do hardware:
    \begin{itemize}
        \item o sintetizador expande e interpreta a lógica descrita;
        \item chamadas repetidas podem representar hardware repetido ou lógica compartilhada conforme o contexto;
        \item o resultado final ainda depende de sinais, processos, clocks, resets e atribuições.
    \end{itemize}

    \vspace{0.5em}
    Portanto, subprogramas melhoram organização e reutilização, mas não substituem o raciocínio sobre o circuito.
\end{frame}

\begin{frame}{Mensagem final}
    Nesta aula, funções e procedimentos apareceram como resposta a problemas concretos:
    \begin{itemize}
        \item a repetição do incremento no timer motivou um procedimento;
        \item a tabela do display motivou uma função;
        \item o VGA mostrou procedimentos separando varredura e desenho;
        \item os testbenches mostraram como validar comportamento sem depender imediatamente do hardware.
    \end{itemize}

    \vspace{0.5em}
    O objetivo é escrever VHDL mais claro sem perder de vista o circuito que está sendo descrito.
\end{frame}

\end{document}
